\documentclass[10pt,a4paper]{amsart}
\usepackage[margin=19mm]{geometry}
\usepackage{amsmath,amssymb,mathtools}
\usepackage[scr=rsfs,cal=euler]{mathalfa}
\usepackage{xfrac}
\usepackage[unicode,hidelinks,bookmarks=false]{hyperref}
\newcommand{\ie}{i.e.\ }
\newcommand{\cf}{cf.\ }
\newcommand{\resp}{respectively\ }
\newcommand{\Subsection}{\subsection}
\providecommand{\nobreakdash}{\nobreak\hbox{-}}
\setlength{\emergencystretch}{2em}
\multlinegap0pt
\usepackage{bm}
\usepackage{bbm}
\usepackage{dsfont}
\usepackage{xspace}
\usepackage{cancel}
\usepackage{mathtools}
\usepackage{amsthm}
\makeatletter
\@ifundefined{lemm}{\newtheorem{lemm}{Lemma}}{}
\@ifundefined{propo}{\newtheorem{propo}[lemm]{Proposition}}{}
\makeatother
\newcommand{\bx}{\mathbf{x}}
\title[Asymptotic stability of a finite sum of solitary waves for t]{Asymptotic stability of a finite sum of solitary waves for the Zakharov-Kuznetsov equation}
\author{Didier Pilod, Fr\'ed\'eric Valet}
\date{Source: arXiv:2312.17721v1 (CC BY 4.0)}
\begin{document}
\maketitle

\noindent\emph{Source and licence.} Asymptotic stability of a finite sum of solitary waves for the Zakharov-Kuznetsov equation. Version arXiv:2312.17721v1 (CC BY 4.0), distributed under the Creative Commons Attribution 4.0 International licence, as recorded in the arXiv licence metadata for this exact version and shown on the abstract page https://arxiv.org/abs/2312.17721v1. Full rights notice: licenses/arxiv-2312.17721-CC-BY-4.0.txt.\par\smallskip

\noindent\textit{Editorial scope.} Counted unit: the Lemm whose source label is \texttt{lemm:quadra\_control\_error} in the article (source0.tex lines 716--721, source label \texttt{lemm:quadra\_control\_error}), delivered together with the article's own complete original proof of it (source0.tex lines 723--803, one \texttt{proof} environment of 81 lines). No mathematics is written, repaired or re-derived: statement and proof are the article's own text line for line. Required context retained uncounted: ZK (lines 116--118); propo:choice\_mod\_param (lines 336--349); eps:ortho (lines 343--348); id:Ri (lines 399--403); est:R1R2.1 (lines 405--418); propo:evol\_mod\_param (lines 425--440); eq:boundz (lines 435--439); eq:c\_n (lines 435--439); lemm:energy (lines 488--493); defi:orbital\_stab (lines 533--535); eq:monotonic\_mass (lines 593--595); ineq:R\_psi\_L2 (lines 603--606). Every definition, display and label of those lines is retained unchanged. Omitted from this package and not redistributed: 1--115, 119--335, 349--398, 404--404, 419--424, 440--487, 494--532, 536--592, 596--602, 607--715, 722--722, 804--2034. Nothing inside a retained excerpt is omitted or altered: every retained body is delivered exactly as the article's lines are written, so no omission receipt of any kind accompanies this package. Citations: the retained excerpts carry no citation command at all, so no bibliography entry of the article is transcribed and no reference is redistributed. Reference adaptation: none was needed and none was made; every reference inside the delivered unit resolves inside it, so no unresolved reference or citation is produced. Printed slips kept literal: no printed word, symbol, hypothesis or number of the counted statement or of its proof was altered. Layout and class: the article's own class is \texttt{amsart} with packages including inputenc, a4wide, amsmath, amsfonts, amssymb, amsthm, mathtools, stmaryrd, bbm, subfiles, appendix, tikz, hyperref, comment; the wrapper is the lane's pinned \texttt{amsart} editorial wrapper at 10pt on A4, which restates the theorem-like environments the retained text needs and adds the article's own macro definitions that the retained text needs (\texttt{\textbackslash bx}), with extra wrapper packages (bm, bbm, dsfont, xspace, cancel, mathtools). The delivered numbering is the unit's own. Assets: no retained span contains a figure environment, a graphics-inclusion command, a PDF-page inclusion, a table or a TikZ picture; no image file is copied into this package, the package declares no asset dependency and no image file is redistributed with it. Licence: version arXiv:2312.17721v1 of this article is distributed under the Creative Commons Attribution 4.0 International licence, as recorded in the arXiv licence metadata for this exact version and shown on the abstract page https://arxiv.org/abs/2312.17721v1. A full rights notice is delivered at licenses/arxiv-2312.17721-CC-BY-4.0.txt. Characters: every retained excerpt is pure ASCII, so the delivered engine, XeLaTeX with its default Latin Modern text font, reports no missing character.
	\begin{align}\label{ZK}%\tag{ZK}
		\partial_t u + \partial_x \left( \Delta u + u^2 \right)=0, 
	\end{align}

	\begin{propo}[Choice of modulation parameters]\label{propo:choice_mod_param}
		Let $0<\underline{c}<\bar{c}$. There exist positive constants $A_1=A_1(\underline{c},\bar{c})$, $\alpha_1^\star=\alpha_1^\star(\underline{c},\bar{c})$, and $Z_1^\star =Z_1^\star(\underline{c},\bar{c})$ such that the following is true. For $0<\alpha< \alpha_1^\star$, $Z>Z_1^\star$ and $\underline{c}<c_2^0<c_1^0<\bar{c}$, let $\mathcal{I}$ be a time interval and $u\in \mathcal{C}(\mathcal{I} : H^1(\mathbb{R}^2))$ be a solution of \eqref{ZK} satisfying $u(t) \in \mathcal{U}_{c_1^0, c_2^0, Z, \alpha}$, for all $t \in \mathcal{I}$.
		Then there exists a unique function $\Gamma=\left(z_1,z_2,\omega_1,\omega_2,c_1,c_2\right) \in \mathcal{C}^1\left(\mathcal{I} : \mathbb{R}^4 \times (0,+\infty)^2\right)$ such that, by defining
		\begin{align}\label{defi:R_n}
			R_i(t,\bx) := Q_{c_i(t)} \left( \bx - \left( z_i(t), \omega_i(t)\right) \right), \ i=1,2, \quad \text{and} \quad \epsilon(t) :=  u(t) - \sum_{i=1}^2 R_i(t) ,
		\end{align}
		we have, for any $t \in \mathcal{I}$ and $i\in \{ 1,2\}$,
		\begin{align}
			& \int R_i(t) \epsilon(t) = \int \partial_x  R_i(t) \epsilon(t) = \int \partial_y  R_i(t)\epsilon(t) = 0, \label{eps:ortho} \\
			& \left\| \epsilon (t) \right\|_{H^1} \leq A_1 \alpha, \label{estimate_modulation:eps} \\ & z(t)= z_1(t)- z_2(t) > Z - A_1 \alpha>\frac{3}4Z, \label{estimate_modulation:z} \\
			&|c_i(t)-c_i^0| \le A_1 \alpha, \label{estimate_modulation:c} \\
            &  \frac{14}{15}\underline{c} < c_2(t) \quad \text{and} \quad c_1(t) < \frac{16}{15}\bar{c}. \label{estimate_modulation:c:bis}
		\end{align}
	\end{propo}

		\begin{align} 
			& \int R_i^2=c_i \int Q^2, \label{id:Ri}\\  
			& \int (\partial_x R_i)^2=\int (\partial_y R_i)^2=c_i^2 \int (\partial_xQ)^2,\label{id:dRi} \\
			&\int (\Lambda_i R_i)^2=c_i^{-1} \int (\Lambda Q)^2.\label{Lambda:dRi}
		\end{align}

		\begin{align} 
			& \left| \int R_i R_j \right| \lesssim  \bar{c}e^{-\frac78  \sqrt{\underline{c}} z}, \label{est:R1R2.1}\\
			& \left| \int R_i \Lambda R_j \right| \lesssim  \bar{c}\underline{c}^{-1}e^{-\frac78 \sqrt{\underline{c}} z} , \label{est:R1R2.2}\\  
			& \left| \int \partial_xR_i R_j \right|+\left| \int \partial_yR_i R_j \right| \lesssim  \bar{c}^{\frac32}e^{-\frac78 \sqrt{\underline{c}} z},
			\label{est:R1R2.3} \\ 
			& \left| \int \partial_xR_i \partial_xR_j \right|+\left| \int \partial_xR_i \partial_yR_j \right|+\left| \int \partial_yR_i \partial_yR_j \right| \lesssim  \bar{c}^{\frac52}e^{-\frac78 \sqrt{\underline{c}} z},
			\label{est:R1R2.4} \\ 
			&\left| \int \partial_xR_i \Lambda R_j \right|+\left| \int \partial_yR_i \Lambda R_j \right| 
			\lesssim \bar{c}^{\frac32}\underline{c}^{-1}e^{-\frac78 \sqrt{\underline{c}} z}, \label{est:R1R2.5} \\ 
			&\left| \int \partial_x(R_iR_j) R_i \right|  
			\lesssim \bar{c}^{\frac52}e^{-\frac{14}{15}\sqrt{\underline{c}} z} \label{est:R1R2.6}\\ 
			&\left| \int \partial_x(R_iR_j) \partial_xR_i \right|  
			\lesssim \bar{c}^{3}e^{-\frac{14}{15}\sqrt{\underline{c}} z} . \label{est:R1R2.7}
		\end{align}

	\begin{propo}[Evolution of the modulation parameters]\label{propo:evol_mod_param}
		Let $0<\underline{c}<\bar{c}$ be fixed. There exist positive constants $k_2=k_2(\underline{c},\bar{c})$, $K_2=K_2(\underline{c},\bar{c})$ and $A_2=A_2(\underline{c},\bar{c})$  such that the following is true.   For two velocities $c_1^0$ and $c_2^0$ satisfying $\underline{c}<c_2^0<c_1^0<\bar{c}$ and 
        \begin{equation} \label{def:sigma}
        \sigma=c_1^0-c_2^0<\min\{1, \underline{c}\},
        \end{equation}
        we define
		\begin{align} \label{def:k2K2}
			\alpha^\star_2:= k_2\sigma \quad \text{and} \quad Z^\star_2 := K_2 | \ln \sigma| .
		\end{align}
		Let $0<\alpha<\alpha_2^{\star}$, $Z>Z_2^{\star}$, $\mathcal{I}$ be a time interval containing $0$ and $u \in C(\mathcal{I} : H^1(\mathbb R^2))$ be a solution of \eqref{ZK} satisfying $u(t) \in \mathcal{U}_{c_1^0, c_2^0, Z, \alpha}$ for all $t \in \mathcal{I}$. Then, under the notation of Proposition \ref{propo:choice_mod_param}, it holds, for all $t \in \mathcal{I}$, 
		\begin{align}
			&\sum_{i=1}^2 \left( \left\vert \dot{c}_i(t) \right\vert  + \left\vert \dot{z}_i(t) -c_i(t) \right\vert + \left\vert \dot{\omega}_i(t) \right\vert \right) \nonumber\\ & \quad \quad \quad \quad \leq A_2 \left( \sum_{i=1}^2\left( \int e^{- \frac{1}{4}\sqrt{\underline{c}} \left\vert \bx - \left( z_i,\omega_i \right) \right\vert } \epsilon^2(\bx) d\bx \right)^\frac12 + e^{-\frac12\sqrt{\underline{c}}\left( Z +  \sigma t \right) } \right), \label{eq:ODE} \\
			& z(t)=z_1(t)-z_2(t)  > \frac34\left(Z +\sigma t\right), \label{eq:boundz} \\
			&\sum_{i=1}^2\left\vert c_i(t) - c_i^0 \right\vert \leq \frac1{32} \sigma. \label{eq:c_n}
		\end{align}
	\end{propo}

	\begin{lemm}\label{lemm:energy}
		Let $0<\underline{c}<\bar{c}$. There exists a positive constant $A_3=A_3(\underline{c},\bar{c})$ such that the following holds. In the setting of Propositions \ref{propo:choice_mod_param} and \ref{propo:evol_mod_param}, we have, for any $t \in \mathcal{I}$,
		\begin{align*}
			\left\vert \sum_{i=1}^2 \left( E(R_i(t))- E(R_i(0))\right) + \frac12 \int \left( \vert \nabla \epsilon \vert^2 - 2 R \epsilon^2 \right)(t) \right\vert \leq A_3 \left( \| \epsilon (0) \|_{H^1}^2 + \| \epsilon(t) \|_{H^1}^3 + e^{-\frac12 \sqrt{\underline{c}}Z} \right).
		\end{align*}
	\end{lemm}

		\begin{align}\label{defi:orbital_stab}
		\gamma:=\frac12 \left( c_1^0 + c_2^0\right), \quad \psi_\gamma(x) := \psi \left( \frac{\sqrt{\gamma}}{2} x \right) \quad \text{and} \quad I(t) := \int u^2(t,\bx) \psi_{\gamma} \left( x - m(t)\right) d\bx. 
		\end{align}

			\begin{align}\label{eq:monotonic_mass}
				I(t)- I(0) \leq A_4  e^{-\frac1{16} \sqrt{\underline{c}}(Z+\sigma t))}.
			\end{align}

\begin{align} 
				&\| R_1 \psi_\gamma'\left( \cdot - m(t)\right) \|_{L^\infty} + \| R_2 \psi_\gamma'\left( \cdot - m(t)\right) \|_{L^\infty} \lesssim \bar{c} e^{-\frac18 \sqrt{\underline{c}} \left( Z + \sigma t\right)} \label{est:Rpsi} \\ 
				&\left\| R_1 \left(\psi_\gamma \left( \cdot - m\right) -1 \right) \right\|_{L^2} + \left\| R_2 \psi_\gamma\left( \cdot - m\right) \right\|_{L^2} \lesssim \bar{c}^\frac12 e^{- \frac1{16} \sqrt{\underline{c}}\left( Z + \sigma t\right)}. \label{ineq:R_psi_L2}
			\end{align}

	\begin{lemm}[Quadratic control of the variation of $c_j(t)$]\label{lemm:quadra_control_error}
		There exists a positive constant $A_8=A_8(\underline{c}, \bar{c})$ such that for any $t\in [0, t^\star]$,
		\begin{align} \label{lemm:quadra_control_error.1}
			\sum_{i=1}^2 \left\vert c_i(t) -c_i(0) \right\vert \leq A_8  \left( \| \epsilon(t) \|_{H_1}^2 + \| \epsilon (0) \|_{H^1}^2 + e^{-\frac1{16} \sqrt{\underline{c}} Z}\right).
		\end{align}
	\end{lemm}

	\begin{proof}
		
		\textit{First step: first order approximation of the localized mass.} Recall the definition of $I$ in \eqref{defi:orbital_stab}. We claim that there exists $A_6 = A_6(\underline{c}, \bar{c})>0$ such that 
		\begin{align}\label{development_M1}
			\left\vert I(t) - I(0) - \left( c_1(t) -c_1(0) \right) \int Q^2 - \int \epsilon^2(t)\psi_\gamma(\cdot-m(t)) \right\vert \leq A_6 \left( \| \epsilon(0) \|_{L^2}^2 + e^{-\frac1{16} \sqrt{\underline{c}} Z} \right)
		\end{align}
   and
   \begin{align}\label{development_M2}
			\left\vert \left( c_2(t) -c_2(0) + c_1(t) -c_1(0) \right) \int Q^2 + \int \epsilon^2(t) \right\vert \leq A_6 \left( \| \epsilon (0) \|_{L^2}^2 + e^{- \frac18 \sqrt{\underline{c}} Z} \right).
		\end{align}
		
		We expand the localized mass $I=I(t)$ as
		\begin{align}\label{eq:M1_developed}
			I = I_1 + I_2 + I_3 + I_4, 
		\end{align}
		with
		\begin{align*}
			 I_1 &:= \int R_1^2+\int \epsilon^2 \psi_\gamma (x-m), \\  
             I_2 &:= 2\int R\epsilon \psi_\gamma (x-m), \\ 
             I_3 &:= \int R_1^2 \left( \psi_\gamma(x-m)-1\right), \\
			 I_4 &:= \int (R-R_1)(R+R_1) \psi_\gamma(x -m).
		\end{align*}
		
		We now estimate $I_2$, $I_3$ and $I_4$ separately. By using \eqref{ineq:R_psi_L2} and  the orthogonality relation \eqref{eps:ortho}, we have
		\begin{align*}
			\left\vert I_2 \right\vert 
			& \lesssim \left\vert  \int \epsilon R_1 \left( \psi_\gamma(\cdot -m) -1\right) \right\vert + \left\vert  \int \epsilon R_2 \psi_\gamma(\cdot-m) \right\vert  \lesssim \| \epsilon \|_{L^2} \bar{c}^\frac12 e^{-\frac1{16}\sqrt{\underline{c}}\left( Z+ \sigma t\right)}; \\
			\left\vert I_3 \right\vert & \lesssim \| R_1 \|_{L^2} \left\| R_1 \left( \psi_\gamma \left( \cdot -m_1\right)-1 \right) \right\|_{L^2} \lesssim \bar{c}^{\frac12} e^{- \frac1{16} \sqrt{\underline{c}}\left( Z + \sigma t\right)}; \\
			\left\vert I_4 \right\vert &\leq \left\| R_2 \psi_\gamma \left( \cdot -m_1 \right) \right\|_ {L^2} \| 2R_1 + R_2 \|_{L^2} \lesssim \bar{c}^{\frac12} e^{- \frac1{16} \sqrt{\underline{c}}\left( Z + \sigma t\right)}.
		\end{align*}
		Taking the difference of \eqref{eq:M1_developed} between the times $t$ and $0$, using \eqref{id:Ri} and the previous bounds conclude the proof of \eqref{development_M1}. 
   
        Similarly, the proof of \eqref{development_M2} follows from the conservation of the mass between the times $t$ and $0$, \eqref{eps:ortho}, \eqref{id:Ri} and \eqref{est:R1R2.1}. 

\medskip
\noindent \textit{Second step: control of the variation of the quadratic speed.}
 	Since $E(Q_c) = c^2 E(Q)$, it follows from Lemma \ref{lemm:energy} and \eqref{eq:boundz} that 
		\begin{align}\label{eq:third_step}
			\left\vert E(Q) \sum_{i=1}^2 \left( c_i(t)^2 - c_i(0)^2 \right) + \frac12 \int \left( \vert \nabla \epsilon \vert^2 - 2 R \epsilon^2 \right)(t) \right\vert \leq A_3 \left( \| \epsilon (0) \|_{H^1}^2 + \| \epsilon (t) \|_{H^1}^3 + e^{- \frac12 \sqrt{\underline{c}}Z} \right).
		\end{align}


\medskip
\noindent	\textit{Third step: Abel resummation argument.} Let us define the quantity
\begin{equation} \label{defi:D}
D(t)=(c_1(0) - c_2(0)) \left( c_1(t) - c_1(0) \right) + c_2(0) \left( c_1(t) -c_1(0) + c_2(t) - c_2(0)  \right).
\end{equation}
Then, we claim that there exists a positive constant $A_7=A_7(\underline{c}, \bar{c})$ such that
\begin{align}
	\left\vert D(t)  \right\vert - \frac12 \sum_{i=1}^2 \vert c_i(t) - c_i(0) \vert^2 \le A_7 \left( \| \epsilon(0) \|_{H^1}^2 +\| \epsilon (t) \|_{H^1}^2 + e^{-\frac18 \sqrt{\underline{c} }Z} \right) .\label{eq:resummation}
\end{align}

Indeed, by using a resummation argument, we have
			$D(t)=\sum_{i=1}^2 c_i(0) \left( c_i(t) -c_i(0) \right)$,
which combined with the triangle inequality and \eqref{eq:third_step} implies
\begin{align}
	\left\vert D(t) \right\vert
	& = \frac12 \left\vert \sum_{i=1}^2 \left( c_i(t)^2 - c_i(0)^2 \right) -\sum_{i=1}^2(c_i(t) -c_i(0))^2 \right\vert \label{eq:Abel}\\
	& \leq \frac{A_3}{2 E(Q)} \left( \| \epsilon (0) \|_{H^1}^2 + \| \epsilon(t) \|_{H^1}^3 + e^{-\frac12 \sqrt{\underline{c}} Z}\right) +\frac{(1+\bar{c})}{2E(Q)} \| \epsilon (t) \|_{H^1}^2 + \frac12 \sum_{i=1}^2 \vert c_i(t) - c_i(0) \vert^2 \nonumber
\end{align}
Then, we conclude the proof of \eqref{eq:resummation} by using $\| \epsilon (t) \|_{H^1} \leq \alpha_5^\star<1$.

\medskip
\noindent \textit{Fourth step: proof of \eqref{lemm:quadra_control_error.1}.} Note that  $\sigma \leq 2(c_1(0)-c_2(0))$ (see\eqref{eq:c_n}) and $\underline{c} < c_2^0$. Combining the statement that if $a \leq b$ with $b \geq 0$, then $\vert a \vert \leq -a +2b$ with  \eqref{eq:monotonic_mass}  \eqref{development_M1}, \eqref{development_M2} and \eqref{eq:resummation} yields
		\begin{align*}
			\MoveEqLeft
			\frac12 \left(\sigma \left\vert c_1(t) -c_1(0) \right\vert + \underline{c} \left\vert c_1(t) - c_1(0) + c_2(t) - c_2(0) \right\vert \right)\\
			& \leq \left( c_1(0) -c_2(0) \right) \left\vert c_1(t) - c_1(0) \right\vert + c_2(0) \left\vert c_1(t) - c_1(0) + c_2(t) - c_2(0) \right\vert \\
			& \leq  -D(t) + 4 \bar{c} \frac{A_6}{\int Q^2}  \| \epsilon (0) \|_{H^1}^2 + \bar{c} \left( 2 A_4 + 3 \frac{A_6}{\int Q^2} \right) e^{- \frac1{16} \sqrt{\underline{c}} Z} \\
			& \leq \left( 4 \bar{c} \frac{A_6}{\int Q^2} + A_7 \right) \left( \| \epsilon (t) \|_{H^1}^2 + \| \epsilon (0) \|_{H^1}^2 \right) + \left( 2 A_4 + 4 \frac{A_6}{\int Q^2} + A_7 \right) e^{- \frac1{16} \sqrt{\underline{c}} Z} \\
   & \quad + \frac12 \sum_{i=1}^2 \vert c_i(t) - c_i(0) \vert^2.
		\end{align*}
		Thus, we deduce from the previous inequality, \eqref{eq:c_n} and $\sigma <\underline{c}$ that
		\begin{align*}
			\MoveEqLeft
			\frac14 \sigma \left[ \vert c_1(t)-c_1(0) \vert + \vert c_2(t)- c_2(0) \vert \right] \\
			& \leq \frac12 \left( \sigma \left\vert c_1(t) - c_1(0) \right\vert + \underline{c} \left\vert c_1(t)- c_1(0) + c_2(t) -c_2(0) \right\vert \right) \\
			& \leq 2\left( 4 \bar{c} \frac{A_6}{\int Q^2} +A_7 \right) \left( \| \epsilon (t) \|_{H^1}^2 + \| \epsilon (0) \|_{H^1}^2 \right) + 2 \left( 2 A_4 + 4 \frac{A_6}{\int Q^2} + A_7 \right) e^{- \frac1{16} \sqrt{\underline{c}} Z},
		\end{align*}
		which concludes the proof of Lemma \ref{lemm:quadra_control_error}.
	\end{proof}

\end{document}
